%%%%%%%%%%  scan idx 200  |  folio 187  %%%%%%%%%%
\emph{Preuve.} L'assertion relative aux générateurs résulte aussitôt de ce que tout faisceau quasi-cohérent sur $S$ est limite inductive filtrante de ses sous-faisceaux quasi-cohérents de type fini (EGA I 9.4.9 et IV 1.7.7). Les autres assertions sont évidentes.

\textbf{Remarque 3.1.1.} Le rédacteur ignore si $\Qcoh(S)$ possède encore assez d'injectifs lorsqu'on suppose seulement $S$ quasi-séparé.

\textbf{Lemme 3.2.} Le foncteur $\phi:\Qcoh(S)\longrightarrow\Mod(S)$ a un adjoint à droite
\[
 Q:\Mod(S)\longrightarrow\Qcoh(S),
\]
(appelé cohéreur).

\emph{Preuve.} Si $S$ est affine, l'existence de $Q$ est immédiate : on prend $Q=\widetilde{\Gamma}_S$ (EGA $\mathrm{Err}_{\mathrm I}$ 1.8.7). Soit maintenant $f:U\longrightarrow S$ un ouvert affine de $S$; il est clair que, pour $E\in\Ob\Qcoh(S)$ et $F\in\Ob\Mod(U)$, on a un isomorphisme canonique
\begin{equation}\tag{1}
 \Hom(E,f_*(F))\simeq \Hom(E,f_*Q(F)).
\end{equation}
Soit alors $(f_i:U_i\longrightarrow S)$ un recouvrement fini de $S$ par des ouverts affines, et pour chaque $(i,j)$ soit $(U_{ijk}\longrightarrow U_{ij})$ un recouvrement fini de $U_{ij}=U_i\cap U_j$ par des ouverts affines ($U_{ij}$ est quasi-compact puisque $S$ est quasi-séparé). Soit $F\in\Ob\Mod(S)$. On a une suite exacte
\begin{equation}\tag{2}
 0\longrightarrow F\longrightarrow \coprod_i f_{i*}(F|U_i)
 \longrightarrow \coprod_{ijk}f_{ijk*}(F|U_{ijk}),
\end{equation}
(où $f_{ijk}$ désigne l'inclusion de $U_{ijk}$ dans $S$). Notons que, si $g:V\longrightarrow U$ est une inclusion d'ouverts affines de $S$, on a une flèche naturelle évidente $Q(F|U)\longrightarrow g_*Q(F|V)$. Par suite, on déduit de (2) une flèche
\begin{equation}\tag{3}
 \coprod_i f_{i*}Q(F|U_i)\longrightarrow
 \coprod_{ijk}f_{ijk*}Q(F|U_{ijk}),
\end{equation}
dont nous désignerons le noyau par $QF$. Soit $E\in\Ob\Qcoh(S)$. Appliquant le foncteur $\Hom(E,.)$ à la suite (2), on obtient, compte tenu de (1), un
